\expandafter\ifx\csname CAPETINCLUDEMODE\endcsname\relax
\documentclass[11pt,a4paper]{article}
% preamble.tex — packages et commandes communes a toutes les fiches CAPET SII IE
\input{glyphtounicode}
\pdfgentounicode=1
\usepackage{cmap}
\usepackage[T1]{fontenc}
\usepackage[utf8]{inputenc}
\usepackage[french]{babel}
\usepackage[a4paper, top=1.6cm, bottom=1.6cm, left=1.4cm, right=1.4cm]{geometry}
\setlength{\headheight}{14.5pt}
\usepackage{amsmath, amssymb, amsfonts}
\usepackage{siunitx}
\sisetup{locale=FR, output-decimal-marker={,}, inter-unit-product=\ensuremath{{}\cdot{}}}
\usepackage{array, booktabs, multirow, tabularx}
\usepackage{xcolor}
\usepackage[most, breakable]{tcolorbox}
\usepackage{multicol}
\usepackage{enumitem}
\usepackage{mathtools}
\usepackage{lmodern}
\usepackage{microtype}
\usepackage{fancyhdr}
\usepackage{tikz}
\usetikzlibrary{positioning,arrows.meta}
\usepackage{lastpage}

% ── Couleurs par domaine ──────────────────────────────────────────────
\definecolor{colElec}{RGB}{26, 82, 148}       % Électrotechnique — bleu foncé
\definecolor{colElecLight}{RGB}{210, 225, 245} % fond bleu clair
\definecolor{colAuto}{RGB}{34, 120, 66}        % Automatique — vert
\definecolor{colMeca}{RGB}{185, 90, 20}        % Mécanique — orange
\definecolor{colTherm}{RGB}{160, 40, 40}       % Thermique — rouge
\definecolor{colQT}{RGB}{232, 244, 233}        % fond question type (vert très léger)
\definecolor{colJury}{RGB}{255, 243, 205}      % fond jury (jaune ambre)
\definecolor{colJuryFrame}{RGB}{200, 130, 0}   % bordure jury
\definecolor{colGrayLight}{RGB}{245, 245, 247} % fond gris neutre
\definecolor{colSeparator}{RGB}{180, 180, 180} % gris séparateur

% ── Style tcolorbox de section principale ────────────────────────────
\tcbset{
  ficheSection/.style={
    breakable,
    enhanced,
    colback=#1!7!white,
    colframe=#1,
    colbacktitle=#1,
    coltitle=white,
    fonttitle=\bfseries\normalsize,
    left=4pt, right=4pt, top=3pt, bottom=4pt,
    before skip=6pt, after skip=6pt,
  }
}

% ── Style encadré question type ──────────────────────────────────────
\tcbset{
  questiontype/.style={
    breakable,
    enhanced,
    colback=colQT,
    colframe=colElec!55!black,
    colbacktitle=colElec!55!black,
    coltitle=white,
    fonttitle=\bfseries\small,
    left=4pt, right=4pt, top=3pt, bottom=3pt,
    before skip=5pt, after skip=5pt,
  }
}

% ── Style encadré jury ───────────────────────────────────────────────
\tcbset{
  jury/.style={
    breakable,
    enhanced,
    colback=colJury,
    colframe=colJuryFrame,
    fonttitle=\bfseries\normalsize,
    left=4pt, right=4pt, top=3pt, bottom=3pt,
  }
}

% ── Mise en page ─────────────────────────────────────────────────────
\pagestyle{fancy}
\fancyhf{}
\renewcommand{\headrulewidth}{0.4pt}
\renewcommand{\footrulewidth}{0.3pt}

\setlist[itemize]{leftmargin=1.2em, itemsep=1pt, topsep=2pt, parsep=0pt}
\setlist[enumerate]{leftmargin=1.2em, itemsep=3pt, topsep=2pt, parsep=0pt}

\setlength{\parskip}{2pt}
\setlength{\columnsep}{10pt}

\def\CAPETBEGINFILE{\begin{document}}
\def\CAPETENDFILE{\end{document}}
\else
\def\CAPETBEGINFILE{}
\def\CAPETENDFILE{}
\fi

\fancyhead[L]{\textbf{\textcolor{colAuto}{T26}} - Modelisation des systemes lineaires (automatique)}
\fancyhead[R]{\small\textit{CAPET SII IE - Session 2026}}
\fancyfoot[C]{\small\thepage\,/\,\pageref{LastPage}}

\CAPETBEGINFILE

% -----------------------------------------------------------------------
%  BLOC EN-TeTE
% -----------------------------------------------------------------------

\begin{tcolorbox}[
  enhanced, breakable,
  colback=colAuto,
  colframe=colAuto,
  left=6pt, right=6pt, top=5pt, bottom=5pt,
  before skip=0pt, after skip=4pt
]
  {\color{white}\bfseries\large T26 \quad Modelisation des systemes lineaires (automatique)}\\[2pt]
  {\color{colAuto!30!white}\small
    \textbf{Domaine :} Automatique - Modelisation et analyse des systemes \quad
    \textbf{Presence :} \textbf{9/13} sujets (2013-2026) - theme \emph{fondamental}
  }\\[1pt]
  {\color{colAuto!20!white}\footnotesize
    Sous-themes : transformee de Laplace - FdT 1\textsuperscript{er}/2\textsuperscript{e} ordre - identification graphique - schema-bloc FTBF - modele MCC - diagramme de Bode - representation d'etat
  }
\end{tcolorbox}

% -----------------------------------------------------------------------
%  1. RAPPELS DE COURS
% -----------------------------------------------------------------------

\begin{tcolorbox}[
  ficheSection=colAuto,
  title={1. Rappels de cours},
]

% -- 1.1 Transformee de Laplace ----------------------------------------
\textbf{1.1 Transformee de Laplace - transformees usuelles}

\smallskip
\begin{center}
\renewcommand{\arraystretch}{1.3}
\begin{tabular}{lll}
\toprule
\textbf{$f(t)$} & \textbf{$F(p)$} & \textbf{Remarque} \\
\midrule
echelon $u(t)$ & $1/p$ & \\
Rampe $t\,u(t)$ & $1/p^2$ & \\
Exponentielle $e^{-t/\tau}$ & $\dfrac{1}{1+\tau p}$ & \\
Derivee $f'(t)$ & $p\,F(p)$ & CI nulles (Heaviside) \\
Integrale $\int_0^t f$ & $F(p)/p$ & \\
\bottomrule
\end{tabular}
\end{center}

\medskip\hrule\vspace{4pt}

% -- 1.2 FdT 1er ordre ------------------------------------------------
\textbf{1.2 Systeme du 1\textsuperscript{er} ordre - forme canonique}

\[
  H(p) = \frac{K}{1 + \tau\,p}
\]
Reponse indicielle : $s(t) = K\,E_0\,(1 - e^{-t/\tau})$ \quad;\quad $t_{r5\%} \approx 3\tau$

\textbf{Identification sur l'essai indiciel :}
\begin{itemize}
  \item $K = s(\infty)/E_0$ (valeur finale / echelon)
  \item $\tau$ : temps pour atteindre $63{,}2\,\%$ de $K\,E_0$ (methode tangente)
\end{itemize}

\medskip\hrule\vspace{4pt}

% -- 1.3 FdT 2e ordre -------------------------------------------------
\textbf{1.3 Systeme du 2\textsuperscript{e} ordre - forme canonique}

\[
  H(p) = \frac{K}{1 + \dfrac{2m}{\omega_0}\,p + \dfrac{p^2}{\omega_0^2}}
  \qquad
  D\% = 100\,e^{-m\pi/\sqrt{1-m^2}} \quad (m<1)
\]
\begin{center}
\renewcommand{\arraystretch}{1.2}
\begin{tabular}{lll}
\toprule
$m > 1$ & Sur-amorti & Pas d'oscillations \\
$m = 1$ & Critique & Retour le plus rapide sans depassement \\
$m < 1$ & Sous-amorti & Oscillations et resonance \\
\bottomrule
\end{tabular}
\end{center}

\medskip\hrule\vspace{4pt}

% -- 1.4 Modele de la MCC ---------------------------------------------
\textbf{1.4 Modele de la machine - courant continu (MCC)}

\[
  u(t) = R\,i(t) + L\,\frac{di}{dt} + e(t)
  \qquad;\qquad
  e(t) = K_e\,\Omega(t)
  \qquad;\qquad
  K_e = K_c
\]

En regime permanent : $U = R\,I + K_e\,\Omega$ ; fonction de transfert moteur-mecanisme :
\[
  \frac{\Omega(p)}{U(p)} = \frac{K_m}{1 + \tau_m\,p}
  \qquad K_m = \frac{1}{K_e} \qquad \tau_m = \frac{J\,R}{K_e^2 + f_v\,R}
\]

\medskip\hrule\vspace{4pt}

% -- 1.5 Schema-bloc - boucle fermee ----------------------------------
\textbf{1.5 Schema-bloc - regles de manipulation}

\[
  H_{BF}(p) = \frac{H_D(p)}{1 + H_D(p)\,H_R(p)} \quad \text{(retour negatif)}
  \qquad
  H_{BO}(p) = H_D(p)\,H_R(p)
\]

Erreur statique en poursuite (classe 0) :
\[
  \varepsilon_s = \frac{E}{1 + K_{BO}} \qquad (\text{nulle si integrateur en BO})
\]

\medskip\hrule\vspace{4pt}

% -- 1.6 Representation d'etat -----------------------------------------
\textbf{1.6 Representation d'etat}

\[
  \dot{\mathbf{x}}(t) = A\,\mathbf{x}(t) + B\,\mathbf{u}(t)
  \qquad;\qquad
  \mathbf{y}(t) = C\,\mathbf{x}(t) + D\,\mathbf{u}(t)
\]
Utilisee pour le modele moyen d'un convertisseur elevateur (2025 E1).

\end{tcolorbox}

% -----------------------------------------------------------------------
%  2. FORMULAIRE SYNTHeTIQUE
% -----------------------------------------------------------------------

\begin{tcolorbox}[
  ficheSection=colAuto,
  title={2. Formulaire synthetique},
]
\small
\begin{multicols}{2}
\renewcommand{\arraystretch}{1.45}
\begin{tabular}{@{}p{5cm}ll@{}}
\toprule
\textbf{Formule} & \textbf{Usage} & \textbf{U.} \\
\midrule
\multicolumn{3}{@{}l}{\footnotesize\textit{Transformee de Laplace}} \\
$\mathcal{L}\{f'(t)\} = p\,F(p)$ & CI nulles & - \\
$\mathcal{L}\{e^{-t/\tau}\} = 1/(1+\tau p)$ & Expo & - \\
\midrule
\multicolumn{3}{@{}l}{\footnotesize\textit{FdT 1\textsuperscript{er} ordre}} \\
$H(p) = K/(1+\tau p)$ & Forme canon. & - \\
$t_{r5\%} \approx 3\tau$ & Rapidite & s \\
$K = s(\infty)/E_0$ & Identification & - \\
$\tau$ : - $63{,}2\,\%$ de $K\,E_0$ & Identification & s \\
\bottomrule
\end{tabular}

\columnbreak

\renewcommand{\arraystretch}{1.45}
\begin{tabular}{@{}p{5cm}ll@{}}
\toprule
\textbf{Formule} & \textbf{Usage} & \textbf{U.} \\
\midrule
\multicolumn{3}{@{}l}{\footnotesize\textit{FdT 2\textsuperscript{e} ordre}} \\
$H = K/(1+2mp/\omega_0+p^2/\omega_0^2)$ & Forme canon. & - \\
$D\% = 100\,e^{-m\pi/\sqrt{1-m^2}}$ & Depassement & \% \\
\midrule
\multicolumn{3}{@{}l}{\footnotesize\textit{Boucle fermee}} \\
$H_{BF} = H_D/(1+H_{BO})$ & FTBF & - \\
$\varepsilon_s = E/(1+K_{BO})$ & Erreur statique & - \\
\midrule
\multicolumn{3}{@{}l}{\footnotesize\textit{MCC}} \\
$U = RI + K_e\Omega$ & Regime perm. & V \\
$K_e = K_c$ & [Ves/rad] = [Nem/A] & \\
\bottomrule
\end{tabular}
\end{multicols}
\normalsize
\end{tcolorbox}

% -----------------------------------------------------------------------
%  3. QUESTIONS TYPES
% -----------------------------------------------------------------------

\begin{tcolorbox}[
  ficheSection=colAuto,
  title={3. Questions types - extraites des sujets reels},
]

% -- QT1 -------------------------------------------------------------

% Mini-procedure de resolution type copie concours
\begin{tcolorbox}[questiontype,title={Procedure type copie concours (5 etapes)}]
\begin{enumerate}[leftmargin=1.5em]
  \item \textbf{Qualifier le modele} : hypotheses, domaine de validite, grandeurs d'entree/sortie.
  \item \textbf{Ecrire les lois} : equations de base avec symboles et unites SI.
  \item \textbf{Isoler l'inconnue} : expression litterale, puis application numerique tracee.
  \item \textbf{Controler l'ordre de grandeur} : coherence physique et unite finale.
  \item \textbf{Conclure concours} : phrase courte orientee decision technique.
\end{enumerate}
\end{tcolorbox}

% Visuel de methode : chemin de resolution CAPET
\begin{center}
\begin{tikzpicture}[font=\small, node distance=7mm and 6mm, >=Stealth]
  \node[draw, rounded corners, fill=colElec!8, inner sep=4pt] (n1) {Notions};
  \node[draw, rounded corners, fill=colElec!8, right=of n1, inner sep=4pt] (n2) {Methode};
  \node[draw, rounded corners, fill=colElec!8, right=of n2, inner sep=4pt] (n3) {Application annale};
  \node[draw, rounded corners, fill=colElec!8, right=of n3, inner sep=4pt] (n4) {Verification jury};
  \draw[->, thick, colElec] (n1) -- (n2);
  \draw[->, thick, colElec] (n2) -- (n3);
  \draw[->, thick, colElec] (n3) -- (n4);
\end{tikzpicture}
\end{center}
\begin{tcolorbox}[questiontype,
  title={QT1 - MCC - equation electrique et identification de $K_e$ \hfill\normalfont\footnotesize E1 2026 (A1-A4)}
]
\textbf{enonce.}
MCC d'une porte de bus (R = 6,6 O ; L = 100 mH ; $K_e$ = 1,044 Veserad$^{-1}$). ecrire l'equation electrique de l'induit. Simplifier en regime permanent. Exprimer $R$ - partir d'un essai rotor bloque ($\Omega = 0$). Retrouver $K_e$ - partir des conditions nominales.

\tcbline

\textbf{Corrige.}

equation electrique de l'induit (loi des mailles) :
\[
  u(t) = R\,i(t) + L\,\frac{di}{dt} + e(t) \quad;\quad e(t) = K_e\,\Omega(t)
\]

Regime permanent ($di/dt = 0$) : $U = R\,I + K_e\,\Omega$.

Rotor bloque ($\Omega = 0$, $e = 0$) : $U_{bl} = R\,I_{bl}$ - $R = U_{bl}/I_{bl}$.

$K_e$ - partir du nominal : $K_e = (U_{nom} - R\,I_{nom})/\Omega_{nom}$.
\end{tcolorbox}

\vspace{4pt}

% -- QT2 -------------------------------------------------------------

\begin{tcolorbox}[questiontype,
  title={QT2 - FdT depuis l'essai indiciel - 1\textsuperscript{er} ordre \hfill\normalfont\footnotesize E1 2026 (C4)}
]
\textbf{enonce.}
L'essai indiciel d'un element de la chaene d'asservissement montre une reponse qui part de 0 et atteint une valeur finale. Proposer un ordre. Construire graphiquement $K$ et $\tau$. Exprimer $H_2(p)$ sous forme canonique.

\tcbline

\textbf{Corrige.}

\textbf{Ordre :} la reponse exponentielle croissante sans oscillation - systeme du \textbf{1\textsuperscript{er} ordre}.

\textbf{Lecture graphique :}
\begin{itemize}
  \item $K = s(\infty)/E_0$ : lire la valeur finale divisee par l'echelon.
  \item $\tau$ : abscisse du point - $63{,}2\,\%$ de $s(\infty)$ (ou intersection de la tangente - l'origine avec le plateau).
\end{itemize}

\textbf{Forme canonique :} $H_2(p) = \dfrac{K}{1 + \tau\,p}$ avec les valeurs lues numeriquement.
\end{tcolorbox}

\vspace{4pt}

% -- QT3 -------------------------------------------------------------

\begin{tcolorbox}[questiontype,
  title={QT3 - FTBF du 2\textsuperscript{e} ordre - depassement nul \hfill\normalfont\footnotesize E1 2026 (C7-C8)}
]
\textbf{enonce.}
La FTBF d'un asservissement avec correcteur $K_C$ est :
$H_{BF}(p) = \dfrac{K_C\,K_m}{1 + (R + K_C\,K_m\,K_{gt})/(K_e)\,\tau_m\,p + \tau_m\,p^2}$.
Montrer qu'elle se met sous la forme canonique du 2\textsuperscript{e} ordre. Exprimer $m$ en fonction de $K_C$. En deduire $K_C$ assurant $D\% = 0$.

\tcbline

\textbf{Corrige.}

Identification par lecture de la forme canonique :
\[
  \omega_0^2 = \frac{K_C\,K_m}{\tau_m} \quad;\quad 2m\,\omega_0 = \frac{R + K_C\,K_m\,K_{gt}}{K_e\,\tau_m}
\]

$D\% = 0$ - $m \geq 1$. Resoudre $m(K_C) = 1$ - valeur limite $K_{Clim}$.

Pour $K_C \leq K_{Clim}$ : $m \geq 1$ - depassement nul. Verifier aussi $M_\varphi$ (Bode).
\end{tcolorbox}

\end{tcolorbox}

% -----------------------------------------------------------------------
%  4. ATTENTES DU JURY
% -----------------------------------------------------------------------

\begin{tcolorbox}[
  jury,
  title={4. Attentes du jury - Erreurs recurrentes et points de vigilance},
]

\begin{enumerate}[label=\textbf{\arabic*.}, leftmargin=*, itemsep=5pt]

  \item \textbf{Maetriser les conditions de Heaviside.} Tout passage equation differentielle - FdT suppose des conditions initiales nulles. Le mentionner explicitement avant d'appliquer la transformee de Laplace.

  \item \textbf{Mettre sous forme canonique avec demonstration.} Recopier la formule finale sans demonstration est sanctionne. Derouler chaque etape d'identification.

  \item \textbf{Identification - $63{,}2\,\%$, pas - $50\,\%$ ni - $90\,\%$.} L'erreur sur la lecture de $\tau$ est frequente - tracer la droite horizontale - $0{,}632 \times K\,E_0$ sur le document reponse.

  \item \textbf{Homogeneite et unites.} $K_e$ [Veserad$^{-1}$] = $K_c$ [NemeA$^{-1}$] numeriquement - ne pas les confondre. Les unites du gain statique $K$ dependent de celles de sortie/entree.

  \item \textbf{Signe du retour dans $H_{BF}$.} $H_{BF} = H_D/(1 + H_{BO})$ : le signe $+$ au denominateur correspond - un retour \textbf{negatif}. Un retour positif donnerait $1 - H_{BO}$ - instabilite.

  \item \textbf{Theme le plus decrie.} Jury 2023 : "L'automatique est un champ disciplinaire auquel les candidats ne peuvent se soustraire. Le jury invite les candidats - revoir les concepts de base." - reviser en priorite.

\end{enumerate}

\medskip
\begin{tcolorbox}[
  colback=colJuryFrame!12!colJury,
  colframe=colJuryFrame,
  left=4pt, right=4pt, top=2pt, bottom=2pt,
  fontupper=\small\itshape
]
Rapport jury 2025 : "Le jury deplore le manque de maetrise des lois basiques de l'electricite necessaires pour etablir les fonctions de transfert. Les expressions litterales homogenes et les unites correctes sont imperatives."\\
Rapport jury 2024 : "La mise sous forme canonique du 2\textsuperscript{e} ordre est exigee avec $m$ et $\omega_0$ explicitement identifies. Ne pas se contenter de factoriser sans nommer les parametres."
\end{tcolorbox}

\end{tcolorbox}

% -----------------------------------------------------------------------
%  PIED - PReSENCE DANS LES SUJETS
% -----------------------------------------------------------------------

\begin{tcolorbox}[
  enhanced, breakable,
  colback=colGrayLight,
  colframe=colSeparator,
  left=4pt, right=4pt, top=3pt, bottom=3pt,
  before skip=4pt, after skip=0pt,
  fonttitle=\small\bfseries,
  title={Presence dans les sujets analyses (2013-2026)},
]
\begin{center}
\footnotesize\renewcommand{\arraystretch}{1.15}
\begin{tabular}{lccccccccccccc}
\toprule
\textbf{Annee} & 13 & 14 & 15 & 16 & 17 & 18 & 19 & 20 & 21 & 22 & 23 & 24 & 25 \\
\midrule
E1 (commune)   & - & - & - & $\checkmark$ & - & $\checkmark$ & $\checkmark$ & - & $\checkmark$ & - & $\checkmark$ & - & $\checkmark$ \\
E2 discipl. IE & - & $\checkmark$ & - & - & - & - & - & $\checkmark$ & - & - & $\checkmark$ & $\checkmark$ & $\checkmark$ \\
E2 appl. comm. & - & - & - & - & - & - & - & - & - & - & - & $\checkmark$ & - \\
\bottomrule
\end{tabular}
\end{center}
\smallskip
\footnotesize\textbf{Sous-themes frequents :}
equation electrique MCC + regime permanent - FdT 1\textsuperscript{er} ordre (identification graphique $K$, $\tau$) - FTBF 2\textsuperscript{e} ordre ($m$, $\omega_0$) - depassement nul - erreur statique (classe 0 vs 1)

\medskip
\textbf{Ordres de grandeur (reperes rapides).}
\begin{itemize}[leftmargin=1.3em]
  \item Reseau BT FR : tension phase-neutre typique \SI{230}{\volt}, tension phase-phase typique \SI{400}{\volt}, frequence \SI{50}{\hertz}.
  \item Courants d'etude frequents en sujet : de quelques \si{\ampere} a quelques dizaines de \si{\ampere}.
  \item En conversion d'energie, un rendement attendu reste generalement entre 0{,}8 et 0{,}98 selon la technologie.
\end{itemize}

\medskip
\textbf{Synthese express (a refaire sans support).}
\begin{enumerate}[leftmargin=1.4em]
  \item Poser les hypotheses du modele et les conventions de signe.
  \item Ecrire 2 a 4 relations cles avec unites.
  \item Realiser un calcul numerique complet et verifier l'ordre de grandeur.
  \item Justifier le resultat par un argument technique de niveau CAPET.
\end{enumerate}

\end{tcolorbox}

\CAPETENDFILE









